\documentclass[10pt]{article}

% Specify the margins and text size.
\setlength{\textwidth}{7in}
\setlength{\textheight}{9.5in}
\setlength{\oddsidemargin}{0pt}
\setlength{\evensidemargin}{0pt}
\setlength{\topmargin}{0pt}
\setlength{\hoffset}{-.3in}
\setlength{\voffset}{-1in}

\setlength{\parskip}{5pt}
\setlength{\parindent}{0pt}

% Load some fonts and symbol packages
\usepackage{latexsym}
\usepackage{pifont}       % contains 'star' symbol for counterinsurgency handbook title
\usepackage{yfonts} 
\usepackage{amsmath}
\usepackage{amsfonts}

\usepackage{graphicx}     % actually, this is loaded by pstricks
\usepackage[T1]{fontenc}
\usepackage{ifthen}
%\usepackage{pstricks,pst-grad,pst-text,pst-node,multido,pst-plot,calc,pst-3dplot}
\usepackage{pstricks,pst-grad,pst-func,pst-text,pst-node,multido,pst-plot,calc,pst-3dplot}
%\usepackage[all]{xy}
%\usepackage{animate}

% The hyperref package inserts links into the pdf.
\definecolor{MyLinkColor}{rgb}{.1,.2,1}
\definecolor{MyCiteColor}{rgb}{.1,1,.2}
\definecolor{MyURLColor}{rgb}{.4,.4,.4}
\usepackage[backref=true,pagebackref=false,hyperindex,colorlinks=true,
  linkcolor=MyLinkColor,urlcolor=MyURLColor]{hyperref}


% The tweaklist package is something I found on the web.  It provides a simple interface
% for making changes to spacing used in the itemize and enumerate environments.  Comment
% this out if you don't care to use tweaklist.
\usepackage{tweaklist}
\renewcommand{\itemhook}{\setlength{\parskip}{2pt}\setlength{\parsep}%
{1pt}\setlength{\topsep}{0pt}\setlength{\itemsep}{0pt}}

\newcommand{\U}{\underline{\hspace{5pt}}}

\usepackage{listings}
\newcommand{\Z}{\hphantom{0}}

\begin{document}
\pagestyle{empty}
\lstset{language=R, showspaces=false, showstringspaces=false}

\href{http://www.su.edu}{\includegraphics[height=1.75cm]{sulogo.eps}}
\vspace{-1.69cm}

{{\ }\hfill\small
\begin{tabular}{cl}
& Math 207\\
& Introduction to Statistics\\
\end{tabular}
}
\setlength{\baselineskip}{1.05\baselineskip}
\bigskip

\begin{center}
\textbf{\large  Exam I}
\end{center}
No books, notes or calculators  are allowed on this exam.
\bigskip

1. A 1994 study of young children found that those with more body fat tended to have
more \lq\lq controlling\rq\rq\ parents.  The \textit{San Francisco Chronicle}  
concluded that \lq\lq Parents of Fat Kids Should Lighten Up.\rq\rq\ (20 points)

\hspace{10pt} a) Was this an observational study or a randomized, controlled experiment?
Explain.
\vspace{.9in}

\hspace{10pt} b) Did the study find an association between parents' behavior
and children's levels of body fat?  Explain.
\vspace{.7in}

\hspace{10pt} c) If controlling behavior causes obesity, could that explain the association?
\vspace{1in}


\hspace{10pt} d) Suppose there were a gene that caused obesity. Why would that \textit{not} explain
the association? 
\vspace{1in}


2. In a health study, subjects' heights have an average of  64 inches and an sd of 3 inches.
(15 points)

\hspace{10pt} a) Make a rough sketch of the histogram of subjects' heights.
Assume a normal (Gaussian) histogram.
\vspace{1.5in}

\hspace{10pt} b) About what percent of the subjects had heights between
61 and 70 inches?  
\vfill
\eject

(25 points)

3.  a) Find the mean, median and sd of the list -4, 5, 3, 5, 5, 12, 9.
\vspace{2.6in}

{\ }\hphantom{n. }b) Which number(s) in the list are within 1 sd of the mean?
\vspace{2.2in}

{\ }\hphantom{n. }c) If you add 4 to each number in the list, what will the new
mean and sd be?
\vspace{2.2in}

{\ }\hphantom{n. }d) If you multiply each number in the list by 4, what will the new
mean and sd be?

\vfill
\eject


4. A histogram of monthly wages for part-time employees is shown below (densities are marked
in parentheses).  Nobody earned more than \$1,000 per month.  (15 points)

\begin{center}
\begin{pspicture}(0,-0.8)(10,2.3)
\psset{yunit=0.1}
\psframe(0,0)(1,10)\rput[b](0.5,10.1){(10)}
\psframe(1,0)(2,4)\rput[b](1.5,4.1){(4)}
\psframe(5,0)(10,10)\rput[b](7.5,10.1){(10)}
\psaxes[Dy=10](0,0)(10,20)
\rput{90}(-1.1,10){\% per \$100}
\rput(5,-10){Wages (hundreds of dollars)}
\end{pspicture}
\end{center}

\hspace{10pt} a) The block
over the class interval from \$200 to \$500 is missing.
Determine how tall it must be then add it to the 
histogram.
\vspace{3in}

\hspace{10pt} b) About what percent of the employees earned between \$600 and \$800 per month?
\vspace{1.75in}

\hspace{10pt} c) About what percent of the employees earned more than \$200 per month?
\vfill
\eject

5. Make a sketch then estimate the area under the normal curve. (5 points each)\vspace{-5pt}

\begin{center}
  a) to the right of $z=1$\hfill
  b) to the left of $z=-6.5$\hfill
  c) between $z=-1$ and $z=3$\hphantom{\hspace{.2in}}
\end{center}
\vspace{3in}

6. (10 points) In class we discussed four types of variables:  ordered, unordered, 
discrete and continuous:\vspace{-5pt}
\begin{center}
{\footnotesize\begin{pspicture}(-1.5,0)(7.5,2.2)
\psset{linewidth=0.02, nodesep=2pt, xunit=3}
\rput(1,2){\rnode{Variables}{Variables}}
\rput(0,1){\rnode{Qualitative}{Qualitative}}
\rput(2,1){\rnode{Quantitative}{Quantitative}}
\rput(1.5,0){\rnode{Discrete}{Discrete}}
\rput(2.5,0){\rnode{Continuous}{Continuous}}
\rput(-0.5,0){\rnode{Ordered}{Ordered}}
\rput(0.5,0){\rnode{Unordered}{Unordered}}
\ncline{Variables}{Qualitative}
\ncline{Variables}{Quantitative}
\ncline{Quantitative}{Discrete}
\ncline{Quantitative}{Continuous}
\ncline{Qualitative}{Ordered}
\ncline{Qualitative}{Unordered}
\end{pspicture}}
\end{center}

In your own words, describe one variable of each of the four types.  Write down some values
for each of your variables. 


\vfill
\eject

%%%%%%%%%%%%%%%%%%%%%%%%%%%%%%%%%%%%%%%%%% HERE

\href{http://www.su.edu}{\includegraphics[height=1.75cm]{sulogo.eps}}
\vspace{-1.69cm}

{{\ }\hfill\small
\begin{tabular}{cl}
& Math 207\\
& Introduction to Statistics\\
\end{tabular}
}
\bigskip

\begin{center}
\textbf{\large  Exam I}
\end{center}
No books, notes or calculators  are allowed on this exam.
\bigskip

1. A histogram of monthly wages for part-time employees is shown below (densities are marked
in parentheses).  Nobody earned more than \$1,000 per month.  (15 points)

\begin{center}
\begin{pspicture}(0,-0.8)(10,2.3)
\psset{yunit=0.1}
\psframe(0,0)(1,20)\rput[b](0.5,20.1){(20)}
\psframe(1,0)(2,12)\rput[b](1.5,12.1){(12)}
\psframe(5,0)(10,7)\rput[b](7.5,7.1){(7)}
\psaxes[Dy=10](0,0)(10,20)
\rput{90}(-1.1,10){\% per \$100}
\rput(5,-10){Wages (hundreds of dollars)}
\end{pspicture}
\end{center}

\hspace{10pt} a) The block
over the class interval from \$200 to \$500 is missing.
Determine how tall it must be then add it to the 
histogram.
\vspace{2.7in}

\hspace{10pt} b) About what percent of the employees earned between \$700 and \$1000 per month?
\vspace{1.75in}

\hspace{10pt} c) About what percent of the employees earned less than \$500 per month?
\vfill
\eject

2. In a health study, subjects' heights have an average of  65 inches and an sd of 5 inches.
(15 points)

\hspace{10pt} a) Make a rough sketch of the histogram of subjects' heights.
Assume a normal (Gaussian) histogram.
\vspace{1.5in}

\hspace{10pt} b) About what percent of the subjects had heights between
55 and 75 inches?  
\vspace{2in}

3. A 1994 study of young children found that those with more body fat tended to have
more \lq\lq controlling\rq\rq\ parents.  The \textit{San Francisco Chronicle}  
concluded that \lq\lq Parents of Fat Kids Should Lighten Up.\rq\rq\ (20 points)

\hspace{10pt} a) Was this an observational study or a randomized, controlled experiment?
Explain.
\vspace{.9in}

\hspace{10pt} b) Did the study find an association between parents' behavior
and children's levels of body fat?  Explain.
\vspace{.7in}

\hspace{10pt} c) If controlling behavior causes obesity, could that explain the association?
\vspace{1in}


\hspace{10pt} d) Suppose there were a gene that caused obesity. Why would that \textit{not} explain
the association? 


\vfill
\eject

(25 points)

4.  a) Find the mean, median and sd of the list 6, 15, 13, 19, 15, 15, 22.
\vspace{2.6in}

{\ }\hphantom{n. }b) Which number(s) in the list are within 1 sd of the mean?
\vspace{2.2in}

{\ }\hphantom{n. }c) If you add 3 to each number in the list, what will the new
mean and sd be?
\vspace{2.2in}

{\ }\hphantom{n. }d) If you multiply each number in the list by 3, what will the new
mean and sd be?

\vfill
\eject





5. (10 points) In class we discussed four types of variables:  ordered, unordered, 
discrete and continuous:\vspace{-5pt}
\begin{center}
{\footnotesize\begin{pspicture}(-1.5,0)(7.5,2.2)
\psset{linewidth=0.02, nodesep=2pt, xunit=3}
\rput(1,2){\rnode{Variables}{Variables}}
\rput(0,1){\rnode{Qualitative}{Qualitative}}
\rput(2,1){\rnode{Quantitative}{Quantitative}}
\rput(1.5,0){\rnode{Discrete}{Discrete}}
\rput(2.5,0){\rnode{Continuous}{Continuous}}
\rput(-0.5,0){\rnode{Ordered}{Ordered}}
\rput(0.5,0){\rnode{Unordered}{Unordered}}
\ncline{Variables}{Qualitative}
\ncline{Variables}{Quantitative}
\ncline{Quantitative}{Discrete}
\ncline{Quantitative}{Continuous}
\ncline{Qualitative}{Ordered}
\ncline{Qualitative}{Unordered}
\end{pspicture}}
\end{center}

In your own words, describe one variable of each of the four types.  Write down some values
for each of your variables. 
\vspace{4in}

6. Make a sketch then estimate the area under the normal curve. (5 points each)\vspace{-5pt}

\begin{center}
  a) to the left of $z=1$\hfill
  b) to the right of $z=-6.5$\hfill
  c) between $z=-2$ and $z=3$\hphantom{\hspace{.2in}}
\end{center}
\vfill
\eject

The \textbf{mean} of a list of numbers is:
\[\mbox{mean}\;=\; \frac{x_1+x_2+\cdots+x_n}{n}\; = \; \frac{1}{n}\sum_{i=1}^n\,x_i\]
The \textbf{sample standard deviation} is
\[\mbox{sd} = \sqrt{\frac{1}{n-1}\,\sum_{i=1}^n\,\left(x_i - \mbox{mean}_x\right)^2}\]
Convert $x$ to \textbf{standard units} using:
\[z = \frac{x - \mbox{mean}}{sd}\]
\bigskip

\begin{center}
\begin{pspicture}(-3,-2.1)(3,1.1)
\psset{yunit=3,linewidth=0.02}
\psline(-3.5,0)(3.5,0)  
  \psline(-3,0)(-3,-0.04) \rput[t](-3,-0.07){\scriptsize -3\hphantom{-}}
  \psline(-2,0)(-2,-0.04) \rput[t](-2,-0.07){\scriptsize -2\hphantom{-}}
  \psline(-1,0)(-1,-0.04) \rput[t](-1,-0.07){\scriptsize -1\hphantom{-}}
  \psline(0,0)(0,-0.04)   \rput[t](0,-0.07){\scriptsize 0}
  \psline(1,0)(1,-0.04)   \rput[t](1,-0.07){\scriptsize 1}
  \psline(2,0)(2,-0.04)   \rput[t](2,-0.07){\scriptsize 2}
  \psline(3,0)(3,-0.04)   \rput[t](3,-0.07){\scriptsize 3}
  \rput[l](3.6,0){\scriptsize $x$}
\psline(0,0)(0,0.5)
  \psline(-0.12,0.5)(0,0.5)    \rput[r](-0.21,0.5){\scriptsize $0.5$}
  \psline(-0.12,0.25)(0,0.25)  \rput[r](-0.21,0.25){\scriptsize $0.25$}
\psGauss[linecolor=blue,linewidth=0.02,sigma=1,mue=0]{-3}{3}
\pnode(-1,-0.15){A}\pnode(1,-0.15){B}
\psbrace[braceWidth=0.02,braceWidthInner=5pt,braceWidthOuter=5pt](A)(B){\rput{90}(0.25,-0.05){\scriptsize 68\%}}
%
\pnode(-2,-0.15){C}\pnode(2,-0.15){D}
\psbrace[braceWidth=0.02,braceWidthInner=25pt,braceWidthOuter=5pt](C)(D){\rput{90}(0.25,-0.05){\scriptsize 95\%}}
%
\pnode(-3,-0.15){E}\pnode(3,-0.15){F}
\psbrace[braceWidth=0.02,braceWidthInner=45pt,braceWidthOuter=5pt](E)(F){\rput{90}(0.25,-0.1){\scriptsize 99.7\%}}
\end{pspicture}
\end{center}
\vfill

The \textbf{mean} of a list of numbers is:
\[\mbox{mean}\;=\; \frac{x_1+x_2+\cdots+x_n}{n}\; = \; \frac{1}{n}\sum_{i=1}^n\,x_i\]
The \textbf{sample standard deviation} is
\[\mbox{sd} = \sqrt{\frac{1}{n-1}\,\sum_{i=1}^n\,\left(x_i - \mbox{mean}_x\right)^2}\]
Convert $x$ to \textbf{standard units} using:
\[z = \frac{x - \mbox{mean}}{sd}\]
\bigskip

\begin{center}
\begin{pspicture}(-3,-2.1)(3,1.1)
\psset{yunit=3,linewidth=0.02}
\psline(-3.5,0)(3.5,0)  
  \psline(-3,0)(-3,-0.04) \rput[t](-3,-0.07){\scriptsize -3\hphantom{-}}
  \psline(-2,0)(-2,-0.04) \rput[t](-2,-0.07){\scriptsize -2\hphantom{-}}
  \psline(-1,0)(-1,-0.04) \rput[t](-1,-0.07){\scriptsize -1\hphantom{-}}
  \psline(0,0)(0,-0.04)   \rput[t](0,-0.07){\scriptsize 0}
  \psline(1,0)(1,-0.04)   \rput[t](1,-0.07){\scriptsize 1}
  \psline(2,0)(2,-0.04)   \rput[t](2,-0.07){\scriptsize 2}
  \psline(3,0)(3,-0.04)   \rput[t](3,-0.07){\scriptsize 3}
  \rput[l](3.6,0){\scriptsize $x$}
\psline(0,0)(0,0.5)
  \psline(-0.12,0.5)(0,0.5)    \rput[r](-0.21,0.5){\scriptsize $0.5$}
  \psline(-0.12,0.25)(0,0.25)  \rput[r](-0.21,0.25){\scriptsize $0.25$}
\psGauss[linecolor=blue,linewidth=0.02,sigma=1,mue=0]{-3}{3}
\pnode(-1,-0.15){A}\pnode(1,-0.15){B}
\psbrace[braceWidth=0.02,braceWidthInner=5pt,braceWidthOuter=5pt](A)(B){\rput{90}(0.25,-0.05){\scriptsize 68\%}}
%
\pnode(-2,-0.15){C}\pnode(2,-0.15){D}
\psbrace[braceWidth=0.02,braceWidthInner=25pt,braceWidthOuter=5pt](C)(D){\rput{90}(0.25,-0.05){\scriptsize 95\%}}
%
\pnode(-3,-0.15){E}\pnode(3,-0.15){F}
\psbrace[braceWidth=0.02,braceWidthInner=45pt,braceWidthOuter=5pt](E)(F){\rput{90}(0.25,-0.1){\scriptsize 99.7\%}}
\end{pspicture}
\end{center}
\bigskip
\eject

\begin{center}
\begin{pspicture}(0,0)(7,4.8)
\psset{xunit=1.4,yunit=1.2,nodesep=2pt}
\rput(4,4){\rnode{Studies}{Studies}}
\rput(5,3){\rnode{NoControls}{No Controls}}
\rput(3,3){\rnode{Controls}{Controls}}
\rput(4,2){\rnode{Historical}{Historical}}
\rput(2,2){\rnode{Contemporaneous}{Contemporaneous}}
\rput(1,1){\rnode{Controlled}{\begin{tabular}{c}Controlled\\ Experiment\end{tabular}}}
\rput(3,1){\rnode{Observational}{\begin{tabular}{c}Observational\\ Studies\end{tabular}}}
\rput(0,0){\rnode{Randomized}{Randomized}}
\rput(2,0){\rnode{NotRandomized}{Not Randomized}}
\ncline{Studies}{Controls}\ncline{Studies}{NoControls}
\ncline{Controls}{Contemporaneous}\ncline{Controls}{Historical}
\ncline{Contemporaneous}{Controlled}\ncline{Contemporaneous}{Observational}
\ncline{Controlled}{Randomized}\ncline{Controlled}{NotRandomized}
%\psframebox(0,0)(4,4)
\end{pspicture}
\end{center}

\vfill
\begin{center}
\begin{pspicture}(0,0)(7,4.8)
\psset{xunit=1.4,yunit=1.2,nodesep=2pt}
\rput(4,4){\rnode{Studies}{Studies}}
\rput(5,3){\rnode{NoControls}{No Controls}}
\rput(3,3){\rnode{Controls}{Controls}}
\rput(4,2){\rnode{Historical}{Historical}}
\rput(2,2){\rnode{Contemporaneous}{Contemporaneous}}
\rput(1,1){\rnode{Controlled}{\begin{tabular}{c}Controlled\\ Experiment\end{tabular}}}
\rput(3,1){\rnode{Observational}{\begin{tabular}{c}Observational\\ Studies\end{tabular}}}
\rput(0,0){\rnode{Randomized}{Randomized}}
\rput(2,0){\rnode{NotRandomized}{Not Randomized}}
\ncline{Studies}{Controls}\ncline{Studies}{NoControls}
\ncline{Controls}{Contemporaneous}\ncline{Controls}{Historical}
\ncline{Contemporaneous}{Controlled}\ncline{Contemporaneous}{Observational}
\ncline{Controlled}{Randomized}\ncline{Controlled}{NotRandomized}
%\psframebox(0,0)(4,4)
\end{pspicture}
\end{center}

\bigskip
\bigskip

\end{document}

4. A histogram of monthly wages for part-time employees is shown below (densities are marked
in parentheses).  Nobody earned more than \$1,000 per month.  (20 points)

\begin{center}
\begin{pspicture}(0,-0.8)(10,2.3)
\psset{yunit=0.1}
\psframe(0,0)(1,10)\rput[b](0.5,10.1){(10)}
\psframe(1,0)(2,5)\rput[b](1.5,5.1){(5)}
\psframe(5,0)(10,5)\rput[b](7.5,5.1){(5)}
\psaxes[Dy=10](0,0)(10,20)
\rput{90}(-1.1,10){\% per \$100}
\rput(5,-10){Wages (hundreds of dollars)}
\end{pspicture}
\end{center}


5. The table below shows results from two experiments designed to test the Salk vaccine in  1954.
(20 points)
\vspace{-5pt}
\begin{center}
{%\footnotesize
\begin{tabular}{cc}
Randomized controlled          & \\
double-blind experiment        & NFIP study\\
   \begin{tabular}{lcc}\hline
                 & Size      & Rate${}^*$\vphantom{\LARGE Y}\\
      \color{black}Treatment  & \color{black}200,000   & {\color{black}{28}}\\
      \color{black}Control    & \color{black}200,000   & {\color{black}{71}}\\
      No consent & 350,000   & 46\\
   \end{tabular}
&\hspace{20pt}
   \begin{tabular}{lcc}
\hline
                 & Size        &  Rate${}^*$\vphantom{\LARGE Y}\\
      \color{black}Grade 2 (vaccine)\hfill        & \color{black}225,000 & 25\\
      \color{black}Grades 1 and 3 (control)        & \color{black}725,000 & {\color{black}{54}}\\
      Grade 2 (no consent)\hfill     & 125,000 & {\color{black}{44}}\\
   \end{tabular}
\end{tabular}}\vspace{-4pt}
\end{center}
{${}^*$Rate is number of cases per 100,000 subjects.}
\smallskip

\hspace{10pt} a) Describe the treatment and control groups and how they were selected 
in the two experiments.  
\vspace{1.75in}


\hspace{10pt} b) What does \textit{consent} mean in the context of the experiments?
\vspace{1.25in}

\hspace{10pt} c) Which rates demonstrate the vaccine's effectiveness in the two experiments 
(there are two rates to compare for each experiment)?  
\vspace{1.25in}

\hspace{10pt} d) Which experiment gave stronger evidence for the vaccine's effectiveness and why?
\vspace{1.5in}

%\hspace{10pt} e) What does \textit{double-blind} mean?
%\vspace{1in}

\hspace{10pt} e) What impact did the different sizes of the six groups have on your conclusions and why?




